Table~\ref{tab:news-taxonomy} shows substantial variation across event
types. Macro news has the largest mean magnitude, 8.96, and the longest
median half-life, 78.1 minutes. Geopolitical news ranks second in
magnitude at 8.08, but its median half-life is 47.1 minutes. Both groups
also rank highest when magnitudes are multiplied by their fitted
window-specific residual loadings. The model therefore associates
broad macroeconomic information with a particularly large and persistent
forecast adjustment, while the distinction for geopolitical articles
is clearer in magnitude than in duration.

The remaining categories illustrate why magnitude and persistence
should be examined separately. Newsletter Series has a relatively
large mean magnitude of 6.85 but a shorter median half-life of
34.4 minutes. Protocol/Infrastructure and Security/Crime have lower
mean magnitudes, 5.79 and 5.81, yet median half-lives of 56.0 and
52.1 minutes. Opinion/Editorial has the lowest mean magnitude, 5.63,
and the shortest median half-life, 25.0 minutes. Product Adoption
also receives a relatively small magnitude, 5.74, although its
43.2-minute median half-life is close to the sample-wide value.
The temporal ordering is thus not a simple restatement of the
magnitude ordering.

These comparisons concern article-level fitted parameters and should
be distinguished from the volume of coverage. Market Commentary is
the largest category, with 2,287 articles, whereas Macro and
Geopolitical contain 450 and 134. Category composition can differ
across the independently fitted windows, and the labels themselves
are automated annotations: 613 articles fall below the confidence
threshold and are assigned to Other. The results provide an
economically recognizable description of the learned news signals,
while leaving causal explanations of category differences open.
